\section{Liquid Welfare and Price of Anarchy}
\label{sec:appendix_welfare}

\subsection{Liquid Welfare}

Following Gaitonde et al.\ (2023, Definition~2.1) and Deng et al.\ (2021), the liquid welfare of an allocation $x$ for $N$ bidders over $T$ rounds is
\begin{equation}
  W(x) = \sum_{i=1}^{N} W_i(x), \qquad
  W_i(x) = \min\!\Big\{B_i,\; \sum_{t=1}^{T} x_{ti}\, v_{ti}\Big\},
  \label{eq:liquid_welfare_def}
\end{equation}
where $v_{ti}$ is bidder~$i$'s valuation in round~$t$, $B_i$ is bidder~$i$'s budget, and $x_{ti} \in [0,1]$ is the fraction of the item allocated to bidder~$i$ in round~$t$.

In the integer case ($x_{ti} \in \{0,1\}$, at most one winner per round), liquid welfare reduces to
\begin{equation}
  W(x) = \sum_{i=1}^{N} \min\!\Big\{B_i,\; \sum_{t:\, \text{winner}_t = i} v_{ti}\Big\},
  \label{eq:liquid_welfare_integer}
\end{equation}
which is how it is computed in the experimental code.

\subsection{LP Relaxation for Optimal Liquid Welfare}

The optimal fractional liquid welfare is the solution to
\begin{equation}
  W^* = \max_{x} \sum_{i=1}^{N} \min\!\Big\{B_i,\; \sum_{t=1}^{T} x_{ti}\, v_{ti}\Big\}
  \quad \text{s.t.} \quad
  \sum_{i=1}^{N} x_{ti} \leq 1 \;\; \forall t, \quad
  x_{ti} \in [0,1] \;\; \forall t,i.
  \label{eq:welfare_opt}
\end{equation}
The $\min$ operator makes this a non-linear program. We linearise it via the following LP:
\begin{equation}
  \text{LP}^* = \max_{x} \sum_{t=1}^{T} \sum_{i=1}^{N} v_{ti}\, x_{ti}
  \quad \text{s.t.} \quad
  \sum_{i=1}^{N} x_{ti} \leq 1 \;\; \forall t, \quad
  \sum_{t=1}^{T} v_{ti}\, x_{ti} \leq B_i \;\; \forall i, \quad
  x_{ti} \in [0,1] \;\; \forall t,i.
  \label{eq:welfare_lp}
\end{equation}

The two programs have the same optimum: $\mathrm{LP}^* = W^*$. At any LP-feasible point $x$, the budget constraint ensures $\min\{B_i, \sum_t v_{ti}\, x_{ti}\} = \sum_t v_{ti}\, x_{ti}$ for each~$i$, so the LP objective equals the liquid welfare, giving $\mathrm{LP}^* \leq W^*$. Conversely, given any optimal solution $x^*$ to~\eqref{eq:welfare_opt}, scaling down allocations for each over-budget bidder~$i$ to $\tilde{x}_{ti} = x^*_{ti} \cdot B_i / (\sum_t v_{ti}\, x^*_{ti})$ produces an LP-feasible point whose objective equals $W(x^*) = W^*$, giving $\mathrm{LP}^* \geq W^*$.

The LP relaxation allows fractional allocations, so $\mathrm{LP}^*$ is an upper bound on the optimal integer liquid welfare. This makes it a conservative benchmark: any PoA computed against this bound is weakly larger than the true integer PoA.

\subsection{Price of Anarchy}

The Price of Anarchy is defined as the ratio of optimal to realised liquid welfare:
\begin{equation}
  \text{PoA} = \frac{W^*}{W(x^{\text{obs}})},
  \label{eq:poa_def}
\end{equation}
where $x^{\text{obs}}$ is the allocation produced by the bidding algorithm. By construction, $\text{PoA} \geq 1$, with values closer to~1 indicating higher efficiency. The literature defines the worst-case PoA as the supremum of this ratio over all instances and equilibria; our experiments report the per-instance empirical PoA.

\subsection{Theoretical Bounds}

\begin{table}[H]
\centering
\small
\caption{Theoretical Price of Anarchy bounds for autobidding.}
\label{tab:poa_bounds}
\begin{tabular}{lll}
\toprule
\textbf{Setting} & \textbf{PoA bound} & \textbf{Source} \\
\midrule
SPA + budget constraints & $\leq 2$ (tight) & Aggarwal et al.\ (2019) \\
FPA + RoS constraints (no budgets) & $= 1$ & Deng et al.\ (2021) \\
FPA + budget constraints & $\leq 2$ (if $v_{ij} \leq B_i$) & Aggarwal et al.\ (2024) \\
Any core auction + learning dynamics & $\leq 2$ & Gaitonde et al.\ (2023) \\
\bottomrule
\end{tabular}
\end{table}

The following bounds are established in the literature for liquid welfare in autobidding environments (Table~\ref{tab:poa_bounds}). The PoA~$= 1$ result of Deng et al.\ (2021) requires return-on-spend constraints, not budget constraints. Experiments~4a and~4b use cumulative budget constraints, for which the applicable theoretical bound is PoA~$\leq 2$ regardless of auction format.
